\documentclass[a4paper]{article}

\usepackage{graphicx}
\usepackage{amsmath,amssymb}
\usepackage{minted}
\usepackage{multirow}
\usepackage{float}
\usepackage{todonotes}
\usepackage{forest}
\usepackage{xstring}
\usepackage{xcolor}
\usepackage[most]{tcolorbox}
\usepackage{xparse}
\usepackage{natbib}
\usepackage{tikz}

\usetikzlibrary{graphs,graphdrawing}
\usegdlibrary{layered}

% External Graphviz support
\input{/home/donkarlo/Dropbox/repo/nd_ai_project/src/nd_ai/knoweldge_representation/language/formal/computer/programming/tex/latex/code_clips/external_graphviz.tex}

% Semantic tag macros
\input{/home/donkarlo/Dropbox/repo/nd_ai_project/src/nd_ai/knoweldge_representation/language/formal/computer/programming/tex/latex/parts/control_sequence/kind/semtag/semtag.tex}

% Link macros
\input{/home/donkarlo/Dropbox/repo/nd_ai_project/src/nd_ai/knoweldge_representation/language/formal/computer/programming/tex/latex/code_clips/seehere.tex}

\title{Waveform}
\author{}
\date{}

\begin{document}
    \maketitle
    
    
    \section{FitzHugh--Nagumo model}
        
        The FitzHugh--Nagumo model (FHN) describes a prototype of an excitable system, such as a neuron.
        It is a simplified dynamical model of neuronal excitability and action-potential-like behavior.
        
        \begin{equation}
            \dot{v}
            =
            c
            \Bigl(
            v
            -
            \frac{v^{3}}{3}
            +
            w
            +
            I
            \Bigr)
        \end{equation}
        
        \begin{equation}
            \dot{w}
            =
            -
            \frac{1}{c}
            \Bigl(
            v
            -
            a
            +
            bw
            \Bigr)
        \end{equation}
        
        \cite{cebrian-lacasa-2025-six-decades-of-the-fitzhugh-nagumo-model-a-guide-through-its-spatio-temporal-dynamics-and-influence-across-disciplines}
        reviews six decades of research on the FitzHugh--Nagumo model as a simplified but powerful description of excitable dynamics.
        It shows that the model captures not only neuronal excitability and oscillations, but also bistability, traveling waves, synchronization, and spatial pattern formation.
        The paper emphasizes that the model has become useful beyond neuroscience, including cardiac dynamics, biology, electronics, optics, mathematics, and physics.
        Its main contribution is to organize the known dynamical regimes and bifurcations of the model, especially in spatially extended and coupled systems.
    
    
    \section{Frequency Modulated M\"obius model}
        
        \cite{rueda-2021-a-novel-wave-decomposition-for-oscillatory-signals}
        and
        \cite{rodriguez-collado-2021-identification-of-action-potential-waveforms}
        propose a wave decomposition method for oscillatory signals based on Frequency Modulated M\"obius (FMM) waves.
        The method represents complex oscillatory patterns as a sum of interpretable wave components with parameters controlling amplitude, phase, location, and shape.
        For action-potential-like curves, its main value is that it provides a citable mathematical basis for approximating signal waveforms by combining several FMM components.
        
        \subsection{Single FMM wave}
            
            A single FMM wave can be written as follows:
            
            \begin{equation}
                W(t)
                =
                A
                \cos
                \Biggl[
                    \beta
                    +
                    2
                    \arctan
                    \Biggl(
                    \omega
                    \tan
                    \Bigl(
                    \frac{t-\alpha}{2}
                    \Bigr)
                    \Biggr)
                    \Biggr]
            \end{equation}
            
            \begin{itemize}
                \item $W(t)$ is the value of the FMM wave at time $t$.
                
                \item $t$ is time.
                
                \item $A$ is the amplitude parameter. It controls the vertical scale of the waveform.
                
                \item $\alpha$ is the location parameter. It controls the horizontal position of the waveform on the time axis.
                
                \item $\beta$ is the phase parameter. It controls the phase shift of the waveform.
                
                \item $\omega$ is the shape parameter. It controls the sharpness, skewness, and temporal deformation of the waveform.
                
                \item $\cos(\cdot)$ is the cosine function used to generate the oscillatory waveform.
                
                \item $\arctan(\cdot)$ is the inverse tangent function used inside the phase transformation.
                
                \item $\tan(\cdot)$ is the tangent function used to transform the time variable before applying the inverse tangent.
                
                \item $\beta + 2 \arctan\bigl(\omega \tan\bigl(\frac{t-\alpha}{2}\bigr)\bigr)$ is the phase function of the FMM wave.
            \end{itemize}
        
        \subsection{Multi-component FMM waveform}
            
            A complex waveform can be modeled as the sum of several FMM waves.
            The general $m$-component model is:
            
            \begin{equation}
                V(t)
                =
                M
                +
                \sum_{j=1}^{m}
                W_j(t)
            \end{equation}
            
            Equivalently:
            
            \begin{equation}
                V(t)
                =
                M
                +
                \sum_{j=1}^{m}
                A_j
                \cos
                \Biggl[
                    \beta_j
                    +
                    2
                    \arctan
                    \Biggl(
                    \omega_j
                    \tan
                    \Bigl(
                    \frac{t-\alpha_j}{2}
                    \Bigr)
                    \Biggr)
                    \Biggr]
            \end{equation}
            
            Here, $M$ is the baseline or vertical offset of the full signal.
            In the case of an action potential, $V(t)$ represents the membrane voltage at time $t$.
            
            \bibliographystyle{plainnat}
            \bibliography{/home/donkarlo/Dropbox/repo/nd_shared_research/bibliography/bibliography.bib}
\end{document}